\documentclass[12pt]{article}
\usepackage{fullpage}
\usepackage[T1]{fontenc}
\usepackage[utf8]{inputenc}
\usepackage{lmodern}
\usepackage{microtype}
\usepackage{amsmath,amssymb,amsthm}
\usepackage{hyperref}

\newtheorem{theorem}{Theorem}
\newtheorem{lemma}{Lemma}
\newtheorem{proposition}{Proposition}
\newtheorem{corollary}{Corollary}
\newtheorem{definition}{Definition}

\newcommand{\Ord}{\mathrm{Ord}}
\newcommand{\Lam}{\Lambda_{\to}}
\newcommand{\to b}{\to_\beta}
\newcommand{\sem}[1]{[\![#1]\!]}
\newcommand{\D}{\mathcal{D}}
\newcommand{\nat}{\mathbb{N}}
\newcommand{\one}{\mathbf{1}}
\newcommand{\nrm}[1]{\lVert #1\rVert}

\begin{document}

\title{An elementary ordinal assignment proving strong normalization\\ of the simply typed lambda calculus}
\author{}
\date{}
\maketitle

\begin{abstract}
We construct an explicit map $\#\colon \Lam \to \Ord$ from simply typed
$\lambda$-terms to ordinals such that every $\beta$-reduction strictly decreases
$\#$, i.e.
\[
  M \to_\beta N \;\Longrightarrow\; \#M > \#N .
\]
The construction is elementary: it interprets each simply typed term as a
\emph{strictly monotone functional} over a hereditarily monotone type structure
built from $\nat$, and reads off $\#M$ as the rank of this functional in a
well-founded order. Nowhere does the definition presuppose strong
normalization; on the contrary, well-foundedness of the ordinals together with
the strict-decrease property yields strong normalization directly.
\end{abstract}

\section{Setup and statement}

\subsection{The calculus}

Simple types are generated by a single base type $o$ and the arrow:
\[
  A,B ::= o \mid A \to B .
\]
We write $A \to B$ for the function type and associate arrows to the right.
Raw terms are
\[
  M,N ::= x \mid (M\,N) \mid (\lambda x{:}A.\,M),
\]
with $x$ ranging over a countable set of variables. A \emph{typing context}
$\Gamma$ is a finite partial map from variables to types. The typing relation
$\Gamma \vdash M : A$ is the usual one:
\[
  \frac{\Gamma(x)=A}{\Gamma \vdash x : A}
  \qquad
  \frac{\Gamma \vdash M : A\to B \quad \Gamma \vdash N : A}{\Gamma \vdash (M\,N):B}
  \qquad
  \frac{\Gamma,\, x{:}A \vdash M : B}{\Gamma \vdash (\lambda x{:}A.\,M): A\to B}.
\]
$\Lam$ denotes the set of typable terms (each carrying a fixed typing
derivation; equivalently we work with Church-style terms in which every bound
variable is annotated by its type, so types are unique). We identify terms up
to renaming of bound variables ($\alpha$-equivalence). Capture-avoiding
substitution is written $M[x:=N]$.

The single-step $\beta$-reduction $\to_\beta$ is the compatible closure of the
contraction rule
\[
  (\lambda x{:}A.\,M)\,N \;\to_\beta\; M[x:=N];
\]
``compatible'' means it may be applied inside any term context:
\begin{equation}\label{eq:compat}
  \frac{M \to_\beta M'}{M\,N \to_\beta M'\,N}
  \qquad
  \frac{N \to_\beta N'}{M\,N \to_\beta M\,N'}
  \qquad
  \frac{M \to_\beta M'}{\lambda x{:}A.\,M \to_\beta \lambda x{:}A.\,M'} .
\end{equation}
Subject reduction holds: if $\Gamma \vdash M:A$ and $M \to_\beta N$ then
$\Gamma \vdash N:A$; we use this implicitly so that all terms appearing in a
reduction sequence have the same type in the same context.

\subsection{Goal}

We construct $\#\colon \Lam \to \Ord$ with
\begin{equation}\label{eq:goal}
  \forall M,N\in\Lam:\quad M \to_\beta N \;\Longrightarrow\; \#M > \#N,
\end{equation}
in a way that does not use strong normalization. Since the ordinals are
well-founded, \eqref{eq:goal} immediately gives strong normalization (Section
\ref{sec:SN}).

\section{The hereditarily monotone type structure}

For each type $A$ we define a set $\D_A$ of ``functionals'' together with a
strict partial order $<_A$. The members of $\D_A$ for $A=o$ are positive
integers; the members for an arrow type are monotone functions.

\begin{definition}[Domains, strict class, and order]\label{def:domains}
By recursion on the type $A$ we simultaneously define a set $\D_A$, a subset
$\mathrm{SM}_A\subseteq\D_A$ (the \emph{strict functionals}), and strict/weak
relations $<_A,\le_A$ on $\D_A$. The recursion is well founded because every
clause for type $A$ refers only to the data already defined at the strictly
smaller component types.

\smallskip\noindent\textbf{Base type.}
\[
  \D_o=\mathrm{SM}_o=\nat_{\ge1}=\{1,2,3,\dots\},\qquad
  m<_o n :\Leftrightarrow m<n,\quad m\le_o n:\Leftrightarrow m\le n.
\]

\smallskip\noindent\textbf{Arrow type $A\to B$.}
$\D_{A\to B}$ is the set of \emph{monotone} functions $f\colon\D_A\to\D_B$, where
monotone means
\[
  a\le_A a'\ \Longrightarrow\ f(a)\le_B f(a')\qquad(a,a'\in\D_A).
\]
The order compares values only at \emph{strict} arguments:
\[
  f<_{A\to B}g :\Leftrightarrow \forall a\in\mathrm{SM}_A:\ f(a)<_B g(a),
  \qquad
  f\le_{A\to B}g :\Leftrightarrow \forall a\in\mathrm{SM}_A:\ f(a)\le_B g(a).
\]
Finally,
\[
  \mathrm{SM}_{A\to B}=\Big\{f\in\D_{A\to B}\ \Big|\
    \begin{array}{l}
      \text{$f$ is strictly monotone on strict arguments:}\\
      \quad a,a'\in\mathrm{SM}_A,\ a<_A a'\ \Rightarrow\ f(a)<_B f(a'),\\
      \text{and } f(a)\in\mathrm{SM}_B \text{ for every } a\in\mathrm{SM}_A
    \end{array}\Big\}.
\]
\end{definition}

\begin{remark}
The relations $<_A,\le_A$ refer only to the (already defined) class
$\mathrm{SM}_A$ at the \emph{argument} type and to $<_B,\le_B$ at the result
type, both strictly smaller than $A\to B$; likewise $\mathrm{SM}_{A\to B}$ uses
only $<_A$ on $\mathrm{SM}_A$ and $<_B,\mathrm{SM}_B$. So the simultaneous
recursion is legitimate. Note that $<_{A\to B}$ is a strict partial order on
$\D_{A\to B}$ in the sense of being irreflexive and transitive (transitivity is
inherited pointwise from $<_B$; irreflexivity holds because $f(a)<_B f(a)$ never
holds, $\mathrm{SM}_A$ being nonempty by Lemma~\ref{lem:unitnorm}); $\le_A$ is a
reflexive transitive preorder, and $f<_A g\Rightarrow f\le_A g$.
\end{remark}

Two operations on functionals will be used repeatedly.

\begin{definition}[Shifting]\label{def:shift}
For $A$ a type, $v\in\D_A$ and $k\in\nat$, the \emph{shift} $v\oplus_A k\in\D_A$
adds $k$ at every base position:
\[
  n\oplus_o k = n+k\ (n\in\D_o);\qquad
  (f\oplus_{A\to B}k)(a) = f(a)\oplus_B k\ \ (a\in\D_A).
\]
\end{definition}

\begin{definition}[Units and norm]\label{def:unit}
By simultaneous recursion on the type (each clause recurses only on strictly
smaller types, so the definition is well founded):
\begin{align*}
  \one_o &= 1, &
  \one_{A\to B}(a) &= \one_B \oplus_B \nrm{a}_A \quad (a\in\D_A);\\
  \nrm{n}_o &= n\ (n\in\D_o), &
  \nrm{f}_{A\to B} &= \nrm{\,f(\one_A)\,}_B \quad (f\in\D_{A\to B}).
\end{align*}
Thus $\one_A\in\D_A$ is a distinguished element and
$\nrm{\cdot}_A\colon\D_A\to\nat$ is a ``norm''.
\end{definition}

We now establish the basic algebraic facts; all are proved by induction on the
type $A$.

\begin{lemma}[Shift is well defined and additive]\label{lem:shift}
For every type $A$, $v\in\D_A$ and $k,k'\in\nat$:
\begin{enumerate}
  \item[(i)] $v\oplus_A k \in\D_A$, and if $v\in\mathrm{SM}_A$ then
    $v\oplus_A k\in\mathrm{SM}_A$;
  \item[(ii)] $(v\oplus_A k)\oplus_A k' = v\oplus_A(k+k')$ and
    $v\oplus_A 0 = v$;
  \item[(iii)] if $k\ge 1$ then $v <_A v\oplus_A k$;
  \item[(iv)] $\oplus_A$ is monotone in its first argument: $u\le_A v$ implies
    $u\oplus_A k \le_A v\oplus_A k$, and $u<_A v$ implies
    $u\oplus_A k <_A v\oplus_A k$.
\end{enumerate}
\end{lemma}

\begin{proof}
\emph{Base $A=o$.} Here $v\oplus_o k = v+k$, and all four statements are
elementary facts about $\nat$: $v+k\in\nat_{\ge1}$; addition is associative with
unit $0$; $k\ge1\Rightarrow v<v+k$; and $u\le v\Rightarrow u+k\le v+k$,
$u<v\Rightarrow u+k<v+k$. Strict monotonicity is vacuous at base type.

\emph{Step $A=C\to D$.} Write $g=v$.
(i) For $a\le_C a'$ we have $g(a)\le_D g(a')$ by monotonicity of $g$, hence
$g(a)\oplus_D k \le_D g(a')\oplus_D k$ by the inductive instance (iv) for $D$;
thus $v\oplus_A k$ is monotone, i.e.\ in $\D_A$. Suppose now
$v=g\in\mathrm{SM}_A$. Then $g$ is strictly monotone, so the same computation
with $a<_C a'$ and inductive (iv) shows $v\oplus_A k$ is strictly monotone; and
for each $a\in\mathrm{SM}_C$, $g(a)\in\mathrm{SM}_D$ (by definition of
$\mathrm{SM}_A$), whence $(v\oplus_A k)(a)=g(a)\oplus_D k\in\mathrm{SM}_D$ by the
inductive instance (i) for $D$. Both requirements of $\mathrm{SM}_A$ hold, so
$v\oplus_A k\in\mathrm{SM}_A$.
(ii) For each $a$, $((g\oplus_D k)\oplus_D k')(a)=(g(a)\oplus_D k)\oplus_D k'
= g(a)\oplus_D(k+k')$ by inductive (ii), and $g\oplus_D 0=g$ likewise;
functions equal pointwise are equal.
(iii) For each $a$, $g(a) <_D g(a)\oplus_D k$ by inductive (iii) (using
$k\ge1$); hence $g <_A g\oplus_A k$.
(iv) For each $a$: if $u\le_A v$ then $u(a)\le_D v(a)$, so
$u(a)\oplus_D k\le_D v(a)\oplus_D k$ by inductive (iv); if $u<_A v$ then
$u(a)<_D v(a)$ for all $a$, so $u(a)\oplus_D k <_D v(a)\oplus_D k$ for all $a$,
i.e.\ $u\oplus_A k <_A v\oplus_A k$.
\end{proof}

\begin{lemma}[Units and norms are strictly monotone]\label{lem:unitnorm}
For every type $A$:
\begin{enumerate}
  \item[(i)] $\one_A \in \mathrm{SM}_A$;
  \item[(ii)] $\nrm{\cdot}_A\colon\D_A\to\nat$ is strictly monotone:
    $u<_A v \Rightarrow \nrm{u}_A < \nrm{v}_A$, and monotone:
    $u\le_A v \Rightarrow \nrm{u}_A \le \nrm{v}_A$;
  \item[(iii)] $\nrm{v\oplus_A k}_A = \nrm{v}_A + k$ for all $v\in\D_A$,
    $k\in\nat$.
\end{enumerate}
\end{lemma}

\begin{proof}
By induction on $A$.

\emph{Base $A=o$.} $\one_o=1$ is (vacuously) strictly monotone;
$\nrm{n}_o=n$ is the identity on $\nat_{\ge1}$, which is strictly monotone and
monotone; and $\nrm{n+k}_o=n+k=\nrm{n}_o+k$.

\emph{Step $A=C\to D$.}
(i) For $a<_C a'$ we have $\nrm{a}_C<\nrm{a'}_C$ by inductive (ii) for $C$,
hence by Lemma~\ref{lem:shift}(iv) for $D$ and the definition of $\one_A$,
\[
  \one_A(a)=\one_D\oplus_D\nrm{a}_C \;<_D\; \one_D\oplus_D\nrm{a'}_C=\one_A(a');
\]
so $\one_A$ is strictly monotone (and monotone, hence in $\D_A$, by the same
computation with $\le$). Moreover, for every $a\in\mathrm{SM}_C$ (indeed for every $a\in\D_C$) the value
$\one_A(a)=\one_D\oplus_D\nrm a_C$ lies in $\mathrm{SM}_D$: by inductive (i)
$\one_D\in\mathrm{SM}_D$, and shifting a strict element by a constant keeps it
strict by Lemma~\ref{lem:shift}(i). Both conditions for $\mathrm{SM}_A$ hold, so
$\one_A\in\mathrm{SM}_A$.
(ii) Suppose $u\le_A v$, i.e.\ $u(a)\le_D v(a)$ for all $a$; taking $a=\one_C$
(which lies in $\D_C$ by inductive (i)) gives $u(\one_C)\le_D v(\one_C)$, so by
inductive (ii) for $D$, $\nrm{u(\one_C)}_D\le\nrm{v(\one_C)}_D$, i.e.
$\nrm{u}_A\le\nrm{v}_A$. If $u<_A v$ then $u(\one_C)<_D v(\one_C)$, and inductive
(ii) gives $\nrm{u}_A<\nrm{v}_A$.
(iii) $\nrm{v\oplus_A k}_A=\nrm{(v\oplus_A k)(\one_C)}_D
=\nrm{v(\one_C)\oplus_D k}_D = \nrm{v(\one_C)}_D+k = \nrm{v}_A+k$, using
Lemma~\ref{lem:shift}(ii) form of the shift on functions and inductive (iii)
for $D$.
\end{proof}

\begin{lemma}[Nonemptiness]\label{lem:nonempty}
Each $\D_A$ is nonempty; indeed $\one_A\in\D_A$.
\end{lemma}
\begin{proof}
Immediate from Lemma~\ref{lem:unitnorm}(i), since strictly monotone elements are
in particular monotone, hence in $\D_A$.
\end{proof}

\section{Interpretation of terms}

\begin{definition}[Valuation, interpretation]\label{def:interp}
Let $\Gamma\vdash M:A$. A \emph{valuation} for $\Gamma$ is a function $\rho$
assigning to each variable $x$ with $\Gamma(x)=C$ an element
$\rho(x)\in\mathrm{SM}_C$. The interpretation $\sem{M}_\rho\in\D_A$ is defined by
recursion on $M$:
\begin{align*}
  \sem{x}_\rho &= \rho(x),\\
  \sem{M\,N}_\rho &= \big(\sem{M}_\rho\big)\!\big(\sem{N}_\rho\big),\\
  \sem{\lambda x{:}C.\,M}_\rho &=
     \Big(\,a \longmapsto \sem{M}_{\rho[x:=a]}\ \oplus_B\ \big(1+\nrm{a}_C\big)\,\Big),
\end{align*}
where in the abstraction clause $\Gamma,x{:}C\vdash M:B$, the bound variable
ranges over $a\in\D_C$, and $\rho[x:=a]$ is $\rho$ updated at $x$ to $a$.
\end{definition}

We must show this is well typed, i.e.\ that each $\sem{M}_\rho$ lies in $\D_A$,
and in fact in $\mathrm{SM}_A$. The crux is that abstraction produces a
\emph{strictly} monotone function even when the bound variable does not occur in
the body, thanks to the additive term $1+\nrm{a}_C$.

\begin{lemma}[Well-definedness and strict monotonicity]\label{lem:wd}
Let $\Gamma\vdash M:A$ and let $\rho$ be a valuation for $\Gamma$. Then
$\sem{M}_\rho\in\mathrm{SM}_A\subseteq\D_A$.
\end{lemma}

\begin{proof}
Throughout this proof we allow \emph{general valuations}: a general valuation for
$\Gamma$ assigns to each $x$ with $\Gamma(x)=C$ an arbitrary element of $\D_C$.
A valuation in the sense of Definition~\ref{def:interp} (assigning $\mathrm{SM}$
values) is called \emph{strict}. The interpretation clauses make sense for
general valuations once each $\sem M_\rho$ is shown to lie in $\D$.

We prove, by induction on the structure of $M$ (with $\Gamma\vdash M:A$):
\begin{itemize}
  \item[(W)] $\sem{M}_\rho\in\D_A$ for every general valuation $\rho$;
  \item[(M)] if $\rho\le\rho'$ pointwise (general valuations) then
    $\sem{M}_\rho\le_A\sem{M}_{\rho'}$;
  \item[(SM)] $\sem{M}_\rho\in\mathrm{SM}_A$ for every \emph{strict} valuation
    $\rho$.
\end{itemize}
The split matters: (W),(M) hold for all $\D$-valued valuations (used to make
abstractions monotone in their possibly-non-strict argument), while (SM) is
asserted only for strict valuations (used when interpreted functionals are
applied to interpreted, hence strict, arguments). The Lemma is the (SM) clause.

\textbf{Case $M=x$, $\Gamma(x)=A$.}
(W): $\sem{x}_\rho=\rho(x)\in\D_A$.
(M): $\rho\le\rho'$ gives $\rho(x)\le_A\rho'(x)$.
(SM): for a strict valuation, $\rho(x)\in\mathrm{SM}_A$.

\textbf{Case $M=P\,Q$, with $\Gamma\vdash P:C\to A$, $\Gamma\vdash Q:C$.}
By induction (W),(M),(SM) hold for $P$ and $Q$.

(W): for general $\rho$, $\sem{P}_\rho\in\D_{C\to A}$ and $\sem{Q}_\rho\in\D_C$,
so $\sem{P\,Q}_\rho=\sem{P}_\rho(\sem{Q}_\rho)\in\D_A$.

(SM): for a strict valuation $\rho$, (SM) for $P,Q$ give
$\sem{P}_\rho\in\mathrm{SM}_{C\to A}$ and $\sem{Q}_\rho\in\mathrm{SM}_C$. By the
definition of $\mathrm{SM}_{C\to A}$ (a strict functional sends strict arguments
to strict values),
\[
  \sem{P\,Q}_\rho=\sem{P}_\rho(\sem{Q}_\rho)\in\mathrm{SM}_A.
\]
No separate ``strictness of values'' argument is needed; that information is
built into the hereditary definition of $\mathrm{SM}$.

(M): for general $\rho\le\rho'$, (M) for $P,Q$ give
$\sem{P}_\rho\le_{C\to A}\sem{P}_{\rho'}$ and $\sem{Q}_\rho\le_C\sem{Q}_{\rho'}$,
hence
\[
  \sem{P\,Q}_\rho
  = \sem{P}_\rho(\sem{Q}_\rho)
   \;\le_A\; \sem{P}_\rho(\sem{Q}_{\rho'})
   \;\le_A\; \sem{P}_{\rho'}(\sem{Q}_{\rho'})
  = \sem{P\,Q}_{\rho'},
\]
the first step by monotonicity of $\sem{P}_\rho\in\D_{C\to A}$ applied to
$\sem{Q}_\rho\le_C\sem{Q}_{\rho'}$, the second by
$\sem{P}_\rho\le_{C\to A}\sem{P}_{\rho'}$ evaluated at $\sem{Q}_{\rho'}$.

\textbf{Case $M=\lambda x{:}C.\,P$ with $\Gamma,x{:}C\vdash P:B$, $A=C\to B$.}
For $a\in\D_C$, $F(a):=\sem{\lambda x.P}_\rho(a)=\sem{P}_{\rho[x:=a]}\oplus_B
(1+\nrm{a}_C)$. Here $\rho[x:=a]$ is a general valuation for $\Gamma,x{:}C$, so
by induction (W) $\sem{P}_{\rho[x:=a]}\in\D_B$, and $F(a)\in\D_B$ by
Lemma~\ref{lem:shift}(i).

\emph{(W): $F$ is monotone, and in fact strictly monotone.} Let $a\le_C a'$.
Then $\rho[x:=a]\le\rho[x:=a']$ as general valuations, so by (M) for $P$,
$\sem{P}_{\rho[x:=a]}\le_B\sem{P}_{\rho[x:=a']}$, whence by
Lemma~\ref{lem:shift}(iv),
\[
  F(a)=\sem{P}_{\rho[x:=a]}\oplus_B(1+\nrm{a}_C)
   \le_B \sem{P}_{\rho[x:=a']}\oplus_B(1+\nrm{a}_C).
\]
If $a<_C a'$, then $\nrm{a}_C<\nrm{a'}_C$ by Lemma~\ref{lem:unitnorm}(ii), so
with $k=\nrm{a'}_C-\nrm{a}_C\ge1$ and Lemma~\ref{lem:shift}(ii)--(iii),
\[
  \sem{P}_{\rho[x:=a']}\oplus_B(1+\nrm{a}_C)
   <_B \sem{P}_{\rho[x:=a']}\oplus_B(1+\nrm{a'}_C)=F(a').
\]
Combining, $a\le_C a'\Rightarrow F(a)\le_B F(a')$ and $a<_C a'\Rightarrow F(a)
<_B F(a')$. Hence $F\in\D_A$ and $F$ is strictly monotone. This is (W); the
strict monotonicity is half of (SM).

\emph{(SM): for a strict valuation $\rho$, $F\in\mathrm{SM}_A$.} We have $F$
strictly monotone. For $a\in\mathrm{SM}_C$, the valuation $\rho[x:=a]$ is strict,
so by (SM) for $P$, $\sem{P}_{\rho[x:=a]}\in\mathrm{SM}_B$, and by
Lemma~\ref{lem:shift}(i) the shift $F(a)\in\mathrm{SM}_B$. Thus $F$ is strictly
monotone and maps $\mathrm{SM}_C$ into $\mathrm{SM}_B$, i.e.\ $F\in
\mathrm{SM}_{C\to B}=\mathrm{SM}_A$.

\emph{(M):} let $\rho\le\rho'$ (general valuations). For every $a\in\D_C$,
$\rho[x:=a]\le\rho'[x:=a]$, so by (M) for $P$ and Lemma~\ref{lem:shift}(iv),
\[
  \sem{\lambda x.P}_\rho(a)=\sem{P}_{\rho[x:=a]}\oplus_B(1+\nrm{a}_C)
  \le_B \sem{P}_{\rho'[x:=a]}\oplus_B(1+\nrm{a}_C)=\sem{\lambda x.P}_{\rho'}(a).
\]
As $a$ was arbitrary, $\sem{\lambda x.P}_\rho\le_A\sem{\lambda x.P}_{\rho'}$.

This completes the induction; the (SM) clause is the assertion of the Lemma.
\end{proof}

\begin{remark}
There is no circularity in the induction: in the application case, the membership
$\sem{P}_\rho\in\mathrm{SM}_{C\to A}$ is simply (SM) for the strictly smaller
term $P$, and the inclusion $\sem{P}_\rho(\sem{Q}_\rho)\in\mathrm{SM}_A$ is then
immediate from the \emph{definition} of $\mathrm{SM}_{C\to A}$. The hereditary
formulation of $\mathrm{SM}$ is precisely what makes the application case a
one-line consequence of the induction hypotheses.
\end{remark}

\section{The substitution and contraction lemmas}

\begin{lemma}[Substitution]\label{lem:subst}
Let $\Gamma,x{:}C\vdash M:A$ and $\Gamma\vdash N:C$, and let $\rho$ be a
general valuation for $\Gamma$ (Lemma~\ref{lem:wd} proof). Then
\[
  \sem{M[x:=N]}_\rho = \sem{M}_{\rho[x:=\sem{N}_\rho]}.
\]
\end{lemma}

\begin{proof}
Note first that $\sem{N}_\rho\in\D_C$ by Lemma~\ref{lem:wd} clause (W), so
$\rho[x:=\sem{N}_\rho]$ is a valuation for $\Gamma,x{:}C$ and the right side is
well defined. We induct on $M$ (using a representative respecting the
$\alpha$-conventions of substitution).

\emph{$M=x$.} $x[x:=N]=N$, so the left side is $\sem{N}_\rho$; the right side is
$\sem{x}_{\rho[x:=\sem N_\rho]}=\sem{N}_\rho$. Equal.

\emph{$M=y\ne x$.} $y[x:=N]=y$, so the left side is $\rho(y)$; the right side is
$\sem{y}_{\rho[x:=\sem N_\rho]}=\rho(y)$ (since $y\ne x$). Equal.

\emph{$M=P\,Q$.} $(P\,Q)[x:=N]=(P[x:=N])\,(Q[x:=N])$, so
\begin{align*}
  \sem{(P\,Q)[x:=N]}_\rho
  &= \sem{P[x:=N]}_\rho\big(\sem{Q[x:=N]}_\rho\big)\\
  &= \sem{P}_{\rho[x:=\sem N_\rho]}\big(\sem{Q}_{\rho[x:=\sem N_\rho]}\big)
   && \text{(induction)}\\
  &= \sem{P\,Q}_{\rho[x:=\sem N_\rho]}.
\end{align*}

\emph{$M=\lambda y{:}D.\,P$ with $y\ne x$ and $y\notin\mathrm{FV}(N)$ (achievable
by $\alpha$-renaming).} Then
$(\lambda y.P)[x:=N]=\lambda y.(P[x:=N])$, and for each $a\in\D_D$,
\begin{align*}
  \sem{\lambda y.(P[x:=N])}_\rho(a)
  &= \sem{P[x:=N]}_{\rho[y:=a]}\oplus_B(1+\nrm{a}_D)\\
  &= \sem{P}_{\rho[y:=a][x:=\sem N_{\rho[y:=a]}]}\oplus_B(1+\nrm{a}_D)
   && \text{(induction).}
\end{align*}
Because $y\notin\mathrm{FV}(N)$ and $y\ne x$, the value $\sem{N}_{\rho[y:=a]}$
does not depend on $a$ and equals $\sem{N}_\rho$ (an interpretation depends only
on the valuation restricted to the free variables of the term; here $y$ is not
free in $N$). Hence the inner valuation is
$\rho[y:=a][x:=\sem N_\rho]=\rho[x:=\sem N_\rho][y:=a]$ (the two updates commute
since $x\ne y$), and the displayed value equals
\[
  \sem{P}_{\rho[x:=\sem N_\rho][y:=a]}\oplus_B(1+\nrm{a}_D)
  = \sem{\lambda y.P}_{\rho[x:=\sem N_\rho]}(a).
\]
As $a\in\D_D$ was arbitrary, the two functions agree, proving the abstraction
case.
\end{proof}

The fact that interpretations depend only on the free variables, used above, is
itself a trivial structural induction and we take it as standard.

\begin{lemma}[Contraction strictly decreases the interpretation]\label{lem:contr}
Let $\Gamma\vdash (\lambda x{:}C.\,M)\,N : B$ (so $\Gamma,x{:}C\vdash M:B$ and
$\Gamma\vdash N:C$) and let $\rho$ be a general valuation for $\Gamma$. Then
\[
  \sem{(\lambda x{:}C.\,M)\,N}_\rho \;>_B\; \sem{M[x:=N]}_\rho .
\]
\end{lemma}

\begin{proof}
Put $a=\sem{N}_\rho\in\D_C$ (Lemma~\ref{lem:wd}, clause (W)). By the abstraction
clause and then Lemma~\ref{lem:subst},
\[
  \sem{(\lambda x.M)\,N}_\rho
  = \sem{\lambda x.M}_\rho(a)
  = \sem{M}_{\rho[x:=a]}\oplus_B\big(1+\nrm{a}_C\big)
  = \sem{M[x:=N]}_\rho \oplus_B \big(1+\nrm{a}_C\big).
\]
Since $1+\nrm{a}_C\ge1$, Lemma~\ref{lem:shift}(iii) gives
$\sem{M[x:=N]}_\rho \oplus_B(1+\nrm{a}_C) >_B \sem{M[x:=N]}_\rho$, which is the
claim.
\end{proof}

\begin{lemma}[Reduction strictly decreases the interpretation]\label{lem:reddec}
Let $\Gamma\vdash M:A$ and $M\to_\beta M'$. Then for every \emph{strict}
valuation $\rho$ for $\Gamma$,
\[
  \sem{M}_\rho \;>_A\; \sem{M'}_\rho .
\]
\end{lemma}

\begin{proof}
By induction on the derivation of $M\to_\beta M'$ from \eqref{eq:compat} and the
contraction rule. Recall (Lemma~\ref{lem:wd}) that under a strict valuation every
subterm interpretation lies in $\mathrm{SM}$ of its type, and that $>_A$ means
$\,f(a)>_B g(a)$ for all $a\in\mathrm{SM}$-arguments when $A$ is an arrow type.

\emph{Base (contraction at the root):}
$M=(\lambda x{:}C.\,P)\,Q\to_\beta P[x:=Q]=M'$. This is exactly
Lemma~\ref{lem:contr} (which holds for general, hence strict, valuations).

\emph{Left of an application:} $M=P\,Q\to_\beta P'\,Q=M'$ with
$P\to_\beta P'$, $\Gamma\vdash P:C\to A$, $\Gamma\vdash Q:C$. By induction
$\sem{P}_\rho>_{C\to A}\sem{P'}_\rho$, i.e.\ $\sem{P}_\rho(c)>_A\sem{P'}_\rho(c)$
for all $c\in\mathrm{SM}_C$. Since $\rho$ is strict, $\sem{Q}_\rho\in\mathrm{SM}_C$
(Lemma~\ref{lem:wd}); taking $c=\sem{Q}_\rho$,
\[
  \sem{P\,Q}_\rho=\sem{P}_\rho(\sem{Q}_\rho)
  >_A \sem{P'}_\rho(\sem{Q}_\rho)=\sem{P'\,Q}_\rho.
\]

\emph{Right of an application:} $M=P\,Q\to_\beta P\,Q'=M'$ with
$Q\to_\beta Q'$, $\Gamma\vdash P:C\to A$. By induction
$\sem{Q}_\rho>_C\sem{Q'}_\rho$, and by Lemma~\ref{lem:wd}
$\sem{Q}_\rho,\sem{Q'}_\rho\in\mathrm{SM}_C$ and
$\sem{P}_\rho\in\mathrm{SM}_{C\to A}$ is strictly monotone on strict arguments.
Hence
\[
  \sem{P\,Q}_\rho=\sem{P}_\rho(\sem{Q}_\rho)
  >_A \sem{P}_\rho(\sem{Q'}_\rho)=\sem{P\,Q'}_\rho.
\]

\emph{Under an abstraction:} $M=\lambda x{:}C.\,P\to_\beta\lambda x{:}C.\,P'=M'$
with $P\to_\beta P'$, $\Gamma,x{:}C\vdash P,P':B$, $A=C\to B$. Let
$a\in\mathrm{SM}_C$ be arbitrary. Then $\rho[x:=a]$ is a \emph{strict} valuation
for $\Gamma,x{:}C$ (all its values, including $a$, are in $\mathrm{SM}$), so by
induction $\sem{P}_{\rho[x:=a]}>_B\sem{P'}_{\rho[x:=a]}$. By
Lemma~\ref{lem:shift}(iv),
\[
  \sem{M}_\rho(a)=\sem{P}_{\rho[x:=a]}\oplus_B(1+\nrm{a}_C)
  >_B \sem{P'}_{\rho[x:=a]}\oplus_B(1+\nrm{a}_C)=\sem{M'}_\rho(a).
\]
As this holds for every $a\in\mathrm{SM}_C$, by the definition of $>_A$ on
$\D_{C\to B}=\D_A$ we conclude $\sem{M}_\rho>_A\sem{M'}_\rho$.
\end{proof}


\section{Well-foundedness and the ordinal assignment}

\begin{lemma}[Well-foundedness of $(\D_A,<_A)$]\label{lem:wf}
For every type $A$, the strict order $<_A$ on $\D_A$ is well founded: there is no
infinite sequence $v_0 >_A v_1 >_A v_2 >_A \cdots$.
\end{lemma}

\begin{proof}
By induction on $A$.

\emph{$A=o$.} $<_o$ is the usual order on $\nat_{\ge1}$, which is well founded.

\emph{$A=C\to B$.} Suppose $v_0>_A v_1>_A\cdots$ were an infinite descending
chain. Since $\one_C\in\mathrm{SM}_C$ (Lemma~\ref{lem:unitnorm}) and $<_A$ quantifies
over strict arguments, we have $v_i(\one_C)>_B v_{i+1}(\one_C)$ for every
$i$. Hence $v_0(\one_C)>_B v_1(\one_C)>_B\cdots$ is an infinite descending chain
in $\D_B$, contradicting the induction hypothesis. So no such chain exists.
\end{proof}

Because $(\D_A,<_A)$ is well founded, the \emph{rank} function
\[
  \mathrm{rk}_A\colon \D_A\to\Ord,\qquad
  \mathrm{rk}_A(v)=\sup\{\,\mathrm{rk}_A(u)+1 : u\in\D_A,\ u<_A v\,\}
\]
is well defined by transfinite recursion on $<_A$, and it is strictly
order-preserving:
\begin{equation}\label{eq:rank}
  u<_A v \;\Longrightarrow\; \mathrm{rk}_A(u) < \mathrm{rk}_A(v).
\end{equation}
(Indeed, if $u<_A v$ then $\mathrm{rk}_A(u)+1\le\mathrm{rk}_A(v)$ by the defining
supremum, so $\mathrm{rk}_A(u)<\mathrm{rk}_A(v)$.)

\begin{definition}[The assignment $\#$]\label{def:sharp}
Fix once and for all, for each typing context $\Gamma$, the canonical valuation
\[
  \rho_0^\Gamma(x) = \one_{\Gamma(x)}\quad (x\in\mathrm{dom}\,\Gamma),
\]
which is a valuation since $\one_C\in\mathrm{SM}_C$ (Lemma~\ref{lem:unitnorm}).
For $M\in\Lam$ with $\Gamma\vdash M:A$ (we use, for each term, the context
consisting of its free variables with their types; this context and type are
determined by $M$), define
\[
  \boxed{\,\#M \;=\; \mathrm{rk}_A\big(\sem{M}_{\rho_0^\Gamma}\big)\in\Ord.\,}
\]
\end{definition}

\begin{theorem}[$\beta$-reduction strictly decreases $\#$]\label{thm:main}
For all $M,N\in\Lam$, if $M\to_\beta N$ then $\#M>\#N$.
\end{theorem}

\begin{proof}
Suppose $\Gamma\vdash M:A$ and $M\to_\beta N$. By subject reduction
$\Gamma\vdash N:A$ in the same context (a $\beta$-step does not create free
variables and may only delete some; if a free variable $x$ of $M$ disappears in
$N$, then $\sem{N}_{\rho_0^\Gamma}$ does not depend on $\rho_0^\Gamma(x)$, so we
may safely use the same valuation $\rho_0^\Gamma$ for both $M$ and $N$, and the
free-variable context of $N$ is a subset of that of $\Gamma$; the value of
$\sem{N}$ is the same whether computed in $\Gamma$ or in the smaller context,
since interpretations depend only on the values at free variables). Applying
Lemma~\ref{lem:reddec} with $\rho=\rho_0^\Gamma$,
\[
  \sem{M}_{\rho_0^\Gamma} \;>_A\; \sem{N}_{\rho_0^\Gamma}.
\]
By \eqref{eq:rank},
\[
  \#M = \mathrm{rk}_A\big(\sem{M}_{\rho_0^\Gamma}\big)
      > \mathrm{rk}_A\big(\sem{N}_{\rho_0^\Gamma}\big) = \#N. \qedhere
\]
\end{proof}

\section{Strong normalization, with no circularity}\label{sec:SN}

\begin{corollary}[Strong normalization]\label{cor:SN}
Every simply typed $\lambda$-term is strongly normalizing: there is no infinite
$\beta$-reduction sequence
\[
  M_0 \to_\beta M_1 \to_\beta M_2 \to_\beta \cdots .
\]
\end{corollary}

\begin{proof}
Suppose such an infinite sequence existed. By subject reduction all $M_i$ have a
common type $A$ in a common context, and by Theorem~\ref{thm:main},
\[
  \#M_0 > \#M_1 > \#M_2 > \cdots
\]
is an infinite strictly descending sequence of ordinals. This contradicts the
well-foundedness of $\Ord$. Hence no infinite reduction sequence exists.
\end{proof}

\begin{remark}[Elementarity]
The map $\#$ is defined purely by recursion on types and terms, using only:
positive integers at the base type; monotone functionals at arrow types; the
pointwise-strict order; and the rank function of a well-founded order. None of
these ingredients refers to the behaviour of reduction, let alone to its
termination. The only fact about reduction we use is the local, syntactic
contraction rule, whose effect on $\#$ is computed in Lemmas
\ref{lem:subst}--\ref{lem:reddec}. Thus Corollary~\ref{cor:SN} is a genuine,
non-circular proof of strong normalization: it deduces termination of reduction
from well-foundedness of the ordinals, rather than assuming termination in the
definition of $\#$.
\end{remark}

\begin{remark}[On the size of $\#$]
The image of $\#$ lands well below $\varepsilon_0$. Indeed $\D_o\cong\omega$, and
an easy induction using Lemma~\ref{lem:wf} shows
$\mathrm{rk}_{A\to B}(v)<\omega^{\,\mathrm{rk}_A(\one_A)\cdot
\mathrm{rk}_B(\,\cdots\,)}$-type bounds; more crudely, every value $\sem M_{\rho_0}$
is a definable functional whose rank is an ordinal below $\omega_n$ (the
finite-level iterated exponential of $\omega$) where $n$ bounds the type level of
$M$. We do not need any such bound for the theorem; transfinite ordinals are
permitted by the problem statement and the bare well-foundedness of $\Ord$ is
all that Corollary~\ref{cor:SN} requires.
\end{remark}

\end{document}
